%% =================================================================
%% Chen et al. Measurement 264 (2026) 120300.
%% Reconstruction of page 13 (printed p. 13).
%% Contents: tail of Section 4.6, full Section 5 Conclusion,
%% CRediT authorship contribution statement, Declaration of competing
%% interest, Acknowledgment, Appendix A (supplementary data),
%% Data availability, References list [1]-[20].
%% =================================================================

\documentclass[11pt]{article}

\usepackage[margin=0.85in]{geometry}
\usepackage{amsmath, amssymb}
\usepackage{mathrsfs}
\usepackage{xcolor}
\usepackage{fancyhdr}
\usepackage{url}

\pagestyle{fancy}
\fancyhf{}
\lhead{\small Z.~Chen et al.}
\rhead{\small Measurement 264 (2026) 120300}
\cfoot{\small 13 $\to$ \arabic{page}}
\renewcommand{\headrulewidth}{0.4pt}

\setcounter{page}{1}
\setcounter{equation}{1}
\setlength{\parskip}{6pt}
\setlength{\parindent}{0pt}

\begin{document}

% ---------- Prose tail of Section 4.6 ----------
polymeric materials typically used for encapsulation. In stark
contrast, the proposed tactile tomography system, which relies on
mechanical sensing rather than acoustic wave propagation, is not
susceptible to such limitations. Its imaging efficacy is inherently
based on the local mechanical interaction (tactile softness)
between the probe and the sample, which effectively bypasses the
interfering effects of the encapsulation layer. This key advantage
allows the tactile tomography system to provide superior imaging
performance for soft-encapsulated structures, highlighting its
unique potential for nondestructive testing in applications where
traditional ultrasonic methods may fail. To demonstrate the
non-destructiveness of the tactile tomography system, a targeted
fatigue test was conducted on the encapsulated FPCB sample. A fixed
point on the encapsulated FPCB was continuously scanned 300 times
using the tactile tomography system under the test conditions of
an applied pressure of 4 MPa, a scanning speed and scanning step
of 0.6 mm/s \& 0.5 mm. Under 300 cycles of constant-pressure
cyclic compression, the compression depth of the probe tip at the
test point remained stable (Fig.~16a), indicating that no stress
distortion occurred in the interaction between the probe and the
encapsulated FPCB during repeated loading, and thus the FPCB did
not suffer cumulative mechanical damage. Fig.~16b shows the
microscope image of the test point of the encapsulated FPCB before
scanning. After 300 cycles of scanning, the post-scanning image is
shown in Fig.~16c, only an almost indiscernible indentation mark
was left on the surface of AB glue encapsulation layer, it neither
cracked nor was punctured, and the underlying FPCB remained intact
with no visible deformation (Fig.~16d). These results fully
demonstrate that the tactile tomography system does not exacerbate
initial or cumulative damage to the encapsulated FPCB.

\section*{5.\ Conclusion}

In this work, a novel tactile tomography system based on mechanical
principles with internal 3D imaging capabilities was developed.
This tactile tomography system can recognize the softness of AB
glues with Shore hardnesses of 3, 7.5, 12.5, 17.5, 25, 30, 35, 50,
and 60 HC. For samples with dimensions less than 12 mm in size and
6 mm in thickness, a novel parameter of ``tactile softness of
materials'' was introduced, enabling the tactile tomography system
to perform internal 3D imaging with a maximum detection depth of 6
mm. The tactile tomography system features a high scanning
resolution of 0.5 mm in both the x and y directions, and 0.1
$\mu$m in the z direction. The tactile tomography system possesses
a scanning accuracy of over 98.9\%, and an optimal scanning speed
and scanning step of 2.4 mm/s \& 0.5 mm. This system is capable of
imaging the internal structure and identifying defects in
encapsulated FPCBs, demonstrating significant potential for
nondestructive testing of flexible wearable devices.

\section*{CRediT authorship contribution statement}

\textbf{Zhiming Chen:} Writing -- review \& editing,
Visualization, Methodology, Investigation.
\textbf{Longxin Qian:} Writing -- original draft, Software, Data
curation.
\textbf{Weihao Liang:} Writing -- original draft, Formal analysis,
Conceptualization.
\textbf{Kaibang Zhang:} Software, Data curation.
\textbf{Zichen Liu:} Investigation, Data curation.
\textbf{Fengming Hu:} Investigation, Data curation.
\textbf{Chaobin Liu:} Investigation, Data curation.
\textbf{Kelan Lu:} Investigation, Data curation.
\textbf{Jingcheng Huang:} Investigation.
\textbf{Haiquan Li:} Methodology.
\textbf{Jinxiu Wen:} Visualization, Methodology, Investigation.
\textbf{Bingpu Zhou:} Methodology, Investigation.
\textbf{Jianyi Luo:} Writing -- review \& editing, Supervision,
Investigation, Conceptualization.

\section*{Declaration of competing interest}

The authors declare that they have no known competing financial
interests or personal relationships that could have appeared to
influence the work reported in this paper.

\section*{Acknowledgment}

This work was financially supported by Innovation and Strong
School Engineering Fund of Guangdong Province (2025KCXTD047),
Guangdong Engineering Technology Research Center (2021J020),
Natural Science Foundation of Guangdong Province (2021A1515011935),
National Natural Science Foundation of China (12004285), Hong Kong
and Macau Joint Research and Development Fund of Wuyi University
(2019WGALH17).

\section*{Appendix A. Supplementary data}

Supplementary data to this article can be found online at\\
\url{https://doi.org/10.1016/j.measurement.2026.120300}.

\section*{Data availability}

Data will be made available on request.

\section*{References}

\begin{footnotesize}
\noindent
[1] M.W.~Vannier, E.V.~Staab, L.P.~Clarke, Matching clinical and
biological needs with emerging imaging technologies, Proc. IEEE 91
(2003) 1562--1573.\\[2pt]
[2] Y.~Zhou, M.~Yuan, J.~Zhang, G.~Ding, S.~Qin, Review of
vision-based defect detection research and its perspectives for
printed circuit board, J. Manuf. Syst. 70 (2023) 557--578.\\[2pt]
[3] C.~Cao, Michael F.~Toney, T.-K.~Sham, R.~Harder, Paul R.~Shearing,
X.~Xiao, J.~Wang, Emerging X-ray imaging technologies for energy
materials, Mater. Today 34 (2020) 132--147.\\[2pt]
[4] V.~Cnudde, M.N.~Boone, High-resolution X-ray computed
tomography in geosciences: a review of the current technology and
applications, Earth Sci. Rev. 123 (2013) 1--17.\\[2pt]
[5] R.~Richards-Kortum, C.~Lorenzoni, Vanderlei S.~Bagnato,
K.~Schmeler, Optical imaging for screening and early cancer
diagnosis in low-resource settings, Nat. Rev. Bioeng. 2 (2023)
25--43.\\[2pt]
[6] A.~Sakdinawat, D.~Attwood, Nanoscale X-ray imaging, Nat.
Photonics 4 (2010) 840--848.\\[2pt]
[7] S.L.~Bridal, J.-M.~Correas, A.~Saied, P.~Laugier, Milestones
on the road to higher resolution, quantitative, and functional
ultrasonic imaging, Proc. IEEE 91 (2003) 1543--1561.\\[2pt]
[8] Elisabeth G.E.~de Vries, L.~Ruijter, Marjolijn N.L.~Hooge,
Rudi A.~Dierckx, Sjoerd G.~Elias, Sjoukje F.~Oosting, Integrating
molecular nuclear imaging in clinical research to improve
anticancer therapy, Nat. Rev. Clin. Oncol. 16 (2019) 241--255.\\[2pt]
[9] T.~Lee, Lili X.~Cai, Victor S.~Lelyveld, A.~Hai, A.~Jasanoff,
Molecular-level functional magnetic resonance imaging of
dopaminergic signaling, Science 344 (2014) 533--535.\\[2pt]
[10] D.E.~Newbury, D.B.~Williams, The electron microscope: the
materials characterization tool of the millennium, Acta Mater. 48
(2000) 323--346.\\[2pt]
[11] A.~Viljoen, M.~Mathelié-Guinlet, A.~Ray, N.~Strohmeyer,
Y.~Oh, P.~Hinterdorfer, Daniel J.~Müller, D.~Alsteens, Yves F.~Dufrêne,
Force spectroscopy of single cells using atomic force microscopy,
Nat. Rev. Methods Primers 1 (2021) 533--535.\\[2pt]
[12] S.~Roath, D.~Newell, A.~Pollack, E.~Alexander, P.-S.~Lin,
Scanning electron microscopy and the surface morphology of human
lymphocytes, Nature 273 (1978) 15--18.\\[2pt]
[13] Y.~Wang, F.~Xie, S.~Ma, L.~Dong, Review of surface profile
measurement techniques based on optical interferometry, Opt. Lasers
Eng. 93 (2017) 164--170.\\[2pt]
[14] J.~Dimastromatteo, T.~Brentnall, Kimberly A.~Kelly, Imaging in
pancreatic disease, Nat. Rev. Gastroenterol. Hepatol. 14 (2017)
97--109.\\[2pt]
[15] E.~Maire, P.J.~Withers, Quantitative X-ray tomography, Int.
Mater. Rev. 59 (2014) 1--43.\\[2pt]
[16] Katharine E.~Thomas, A.~Fotaki, Rene M.~Botnar, Vanessa M.~Ferreira,
Imaging Methods: Magnetic Resonance Imaging, Circulation 122.014068
(2023).\\[2pt]
[17] N.~Li, L.~Wang, J.~Jia, Y.~Yang, A novel method for the image
quality improvement of ultrasonic tomography, IEEE Trans. Instrum.
Meas. 70 (2021) 1--10.\\[2pt]
[18] J.~Kim, W.~Brown, Jason R.~Maher, H.~Levinson, A.~Wax,
Functional optical coherence tomography: principles and progress,
Physics in Medicine and Biology 60 (2015) R211--R237.\\[2pt]
[19] J.~Rong, A.~Haider, Troels E.~Jeppesen, L.~Josephson,
Steven H.~Liang, Radiochemistry for positron emission tomography,
Nat. Commun. 14 (2023) 17200--17206.\\[2pt]
[20] K.~Shin, D.~Kim, H.~Park, M.~Sim, H.~Jang, J.~Sohn, S.~Cha,
J.~Jang, Artificial tactile sensor with pin-type module for depth
profile and surface topography detection,
\end{footnotesize}

\end{document}
